% GENERATED FILE. Edit canonical lessons, then run npm run generate:books.
% canonical-source-hash: fnv1a64:cccf93b02efaa2b8
% canonical-lessons: TE-C118-tamata, TE-C118-bangaladumpa, TE-C118-vankaya, TE-C118-allam, TE-C118-vellulli

\chapter{Tomato, Potato, Aubergine, Ginger, Garlic}
\label{ch:te-tomato-potato-aubergine-ginger-garlic}
\hlchaptermodality{118}

\begin{chapteropening}
\textbf{By the end of this chapter:} \emph{I can say a tomato, a potato, an aubergine, ginger and garlic.}

Complete the last lesson of chapter 118: i can say a tomato, a potato, an aubergine, ginger and garlic.
\end{chapteropening}

\section[\texorpdfstring{\te{టమాటా} (ṭamāṭā)}{టమాటా (ṭamāṭā)}]{\te{టమాటా} (ṭamāṭā) — a tomato}

\label{lesson:TE-C118-tamata}



\begin{quote}
Before the new one: say the Telugu for coconut, then the Telugu for an onion.
\end{quote}

\subsection*{You'll want to know: \te{టమాటా}}
\textbf{\te{టమాటా}} — \emph{ṭamāṭā} — \textquotedblleft{}a tomato\textquotedblright{}.

\begin{grammarlens}[title={What you've built}]
One more word for this chapter.
\end{grammarlens}

\subsection*{Your turn}
\begin{itemize}
\raggedright
  \item \emph{Say it:} \emph{ṭamāṭā}
  \item \emph{Say it:} \emph{ṭamāṭā}, once more
  \item \emph{Say it:} say \emph{ṭamāṭā}
  \item \emph{Recall:} say the Telugu for coconut, then the Telugu for an onion, then say \emph{ṭamāṭā} again
\end{itemize}

\begin{tcolorbox}[breakable,colback=teal!4,colframe=teal!35!black,title={Before you move on}]
What does \te{టమాటా} mean? (\textquotedblleft{}A tomato\textquotedblright{}.) Say it once more.
\end{tcolorbox}
\section[\texorpdfstring{\te{బంగాళాదుంప} (baṅgāḷādumpa)}{బంగాళాదుంప (baṅgāḷādumpa)}]{\te{బంగాళాదుంప} (baṅgāḷādumpa) — a potato}

\label{lesson:TE-C118-bangaladumpa}



\begin{quote}
Before the new one: say the Telugu for an onion, then the Telugu for a tomato.
\end{quote}

\subsection*{You'll want to know: \te{బంగాళాదుంప}}
\textbf{\te{బంగాళాదుంప}} — \emph{baṅgāḷādumpa} — \textquotedblleft{}a potato\textquotedblright{}.

\begin{grammarlens}[title={What you've built}]
One more word for this chapter.
\end{grammarlens}

\subsection*{Your turn}
\begin{itemize}
\raggedright
  \item \emph{Say it:} \emph{baṅgāḷādumpa}
  \item \emph{Say it:} \emph{baṅgāḷādumpa}, once more
  \item \emph{Say it:} say \emph{baṅgāḷādumpa}
  \item \emph{Recall:} say the Telugu for an onion, then the Telugu for a tomato, then say \emph{baṅgāḷādumpa} again
\end{itemize}

\begin{tcolorbox}[breakable,colback=teal!4,colframe=teal!35!black,title={Before you move on}]
What does \te{బంగాళాదుంప} mean? (\textquotedblleft{}A potato\textquotedblright{}.) Say it once more.
\end{tcolorbox}
\section[\texorpdfstring{\te{వంకాయ} (vaṅkāya)}{వంకాయ (vaṅkāya)}]{\te{వంకాయ} (vaṅkāya) — an aubergine}

\label{lesson:TE-C118-vankaya}



\begin{quote}
Before the new one: say the Telugu for a tomato, then the Telugu for a potato.
\end{quote}

\subsection*{You'll want to know: \te{వంకాయ}}
\textbf{\te{వంకాయ}} — \emph{vaṅkāya} — \textquotedblleft{}an aubergine\textquotedblright{}.

\begin{grammarlens}[title={What you've built}]
One more word for this chapter.
\end{grammarlens}

\subsection*{Your turn}
\begin{itemize}
\raggedright
  \item \emph{Say it:} \emph{vaṅkāya}
  \item \emph{Say it:} \emph{vaṅkāya}, once more
  \item \emph{Say it:} say \emph{vaṅkāya}
  \item \emph{Recall:} say the Telugu for a tomato, then the Telugu for a potato, then say \emph{vaṅkāya} again
\end{itemize}

\begin{tcolorbox}[breakable,colback=teal!4,colframe=teal!35!black,title={Before you move on}]
What does \te{వంకాయ} mean? (\textquotedblleft{}An aubergine\textquotedblright{}.) Say it once more.
\end{tcolorbox}
\section[\texorpdfstring{\te{అల్లం} (allaṁ)}{అల్లం (allaṁ)}]{\te{అల్లం} (allaṁ) — ginger}

\label{lesson:TE-C118-allam}



\begin{quote}
Before the new one: say the Telugu for a potato, then the Telugu for an aubergine.
\end{quote}

\subsection*{You'll want to know: \te{అల్లం}}
\textbf{\te{అల్లం}} — \emph{allaṁ} — \textquotedblleft{}ginger\textquotedblright{}.

\begin{grammarlens}[title={What you've built}]
One more word for this chapter.
\end{grammarlens}

\subsection*{Your turn}
\begin{itemize}
\raggedright
  \item \emph{Say it:} \emph{allaṁ}
  \item \emph{Say it:} \emph{allaṁ}, once more
  \item \emph{Say it:} say \emph{allaṁ}
  \item \emph{Recall:} say the Telugu for a potato, then the Telugu for an aubergine, then say \emph{allaṁ} again
\end{itemize}

\begin{tcolorbox}[breakable,colback=teal!4,colframe=teal!35!black,title={Before you move on}]
What does \te{అల్లం} mean? (\textquotedblleft{}Ginger\textquotedblright{}.) Say it once more.
\end{tcolorbox}
\section[\texorpdfstring{\te{వెల్లుల్లి} (vellulli)}{వెల్లుల్లి (vellulli)}]{\te{వెల్లుల్లి} (vellulli) — garlic}

\label{lesson:TE-C118-vellulli}



\begin{quote}
Before the new one: say the Telugu for an aubergine, then the Telugu for ginger.
\end{quote}

\subsection*{You'll want to know: \te{వెల్లుల్లి}}
\textbf{\te{వెల్లుల్లి}} — \emph{vellulli} — \textquotedblleft{}garlic\textquotedblright{}.

\begin{grammarlens}[title={What you've built}]
One more word for this chapter.
\end{grammarlens}

\subsection*{Your turn}
\begin{itemize}
\raggedright
  \item \emph{Say it:} \emph{vellulli}
  \item \emph{Say it:} \emph{vellulli}, once more
  \item \emph{Say it:} say \emph{vellulli}
  \item \emph{Recall:} say the Telugu for an aubergine, then the Telugu for ginger, then say \emph{vellulli} again
\end{itemize}

\begin{tcolorbox}[breakable,colback=teal!4,colframe=teal!35!black,title={Before you move on}]
What does \te{వెల్లుల్లి} mean? (\textquotedblleft{}Garlic\textquotedblright{}.) Say it once more.
\end{tcolorbox}
